\chapter{Entrer sur un nouveau marché}\label{ch:fm:entering-a-new-market}

Le plan annonce huit mois avant la première transaction. La convention de compensation seule en prend cinq, et elle ne peut commencer que lorsque l'entité juridique existe et
dispose de son identifiant~; la bourse ne certifiera le logiciel qu'une fois que la société de compensation aura fixé les limites de la société~; et l'enregistrement réglementaire
qu'exige l'adhésion prend à lui seul quatre mois, ne démarrant qu'une fois l'entité constituée. Représentée sous forme de réseau, l'entrée dure dix mois, et le
goulot d'étranglement n'est pas la convention de compensation que tout le monde surveillait, mais l'enregistrement que personne n'avait inscrit sur le graphique.

\section{Pourquoi entrer~: l'edge et la taille du gain}

Une société de trading entre sur un nouveau marché (une classe d'actifs, une région, une plateforme de négociation) parce qu'un edge qu'elle possède ailleurs devrait y fonctionner, ou parce qu'un client ou
une contrepartie qu'elle sert déjà y opère. Le rapport annuel 2025 de Flow Traders est un exemple public du second type de raisonnement~: il décrit
l'expansion de sa présence en Asie et «~le lancement d'une activité de négociation active en Chine, un marché de croissance attractif, en phase avec nos atouts dans les ETP~», avec
la valeur négociée de ses produits négociés en bourse (ETP) en Asie passant de 114~milliards d'euros en 2024 à 152~milliards d'euros en 2025. Quelle que soit la raison, l'entrée coûte
de l'argent pendant des mois avant la première transaction, et le gain reste incertain tant que la société n'a pas traité pendant un certain temps.

\begin{definition}[Plan d'entrée, délai jusqu'à la première transaction]\label{def:fm:entering-a-new-market:plan}
Un \emph{plan d'entrée}\index{plan d'entrée} recense les chantiers nécessaires pour négocier sur un nouveau marché (entité juridique, \omterm{def:fm:regulators-and-licences:perimeter}{agrément}, compensation,
adhésion, connectivité, données, personnel, logiciel, certification), leurs durées, leurs coûts et leurs dépendances, la montée en régime attendue des revenus et les
points où la société décide de poursuivre ou non. Le \emph{délai jusqu'à la première transaction}\index{délai jusqu'à la première transaction} est le temps écoulé entre la décision
d'entrer et la première transaction réelle, fixé par la plus longue chaîne de chantiers dépendants.
\end{definition}

\section{La liste de contrôle~: accès, compensation, données, juridique, effectifs}

Les chantiers sont les chapitres de cette série. L'entité a besoin d'un identifiant d'entité juridique (livre~2) avant que la plupart des contreparties n'ouvrent un
compte~; l'\omterm{def:fm:regulators-and-licences:perimeter}{agrément} ou l'enregistrement suit la carte du chapitre~17 et ses délais réglementaires~; la compensation passant par un futures commission merchant (FCM) ou par
un membre compensateur (livre~1) exige des contrôles KYC, des conventions juridiques (chapitre~18) et des limites~; l'adhésion et la certification (livre~13) suivent~;
la connectivité (livre~14), les licences de données de marché (chapitre~22) et le personnel (chapitre~10) avancent en parallèle. L'exemple du chapitre est un marché à terme dans une nouvelle
région (le \cref{tab:fm:entering-a-new-market:tasks})~; les durées et les coûts sont des données illustratives.

\begin{table}[ht]
\centering\small
\begin{tabular}{lrrlr}
\toprule
chantier & mois & k\$/mois & après & marge \\
\midrule
entité juridique & 2.0 & 25 & -- & 0 \\
identifiant d'entité juridique & 0.5 & 2 & entité & 1.0 \\
enregistrement réglementaire & 4.0 & 30 & entité & 0 \\
convention de compensation & 5.0 & 20 & identifiant & 1.0 \\
limites de compensation fixées & 0.5 & 5 & convention de compensation & 1.0 \\
adhésion à la bourse & 3.0 & 15 & enregistrement & 0 \\
connectivité & 3.0 & 40 & entité & 4.0 \\
licences de données de marché & 2.0 & 10 & entité & 6.0 \\
embauches & 4.0 & 60 & -- & 2.0 \\
adaptation du logiciel & 3.0 & 80 & embauches & 2.0 \\
certification & 1.0 & 10 & limites, adhésion, & 0 \\
 & & & connectivité, logiciel & \\
\bottomrule
\end{tabular}
\caption{Les chantiers de l'entrée, avec leur marge en mois dans le calendrier au plus tôt (illustratif). Données~: \texttt{fm\_entry.TASKS}, \texttt{firm.entryplan.schedule}.}\label{tab:fm:entering-a-new-market:tasks}
\end{table}

\section{Le chemin critique}

\begin{definition}[Méthode du chemin critique]\label{def:fm:entering-a-new-market:cpm}
La \emph{méthode du chemin critique}\index{méthode du chemin critique} ordonnance un projet donné sous forme de tâches assorties de durées et de contraintes de précédence~: chaque tâche
commence lorsque tous ses prédécesseurs sont achevés~; la durée du projet est celle du plus long chemin du réseau~; la marge d'une tâche est le retard
qu'elle peut subir sans retarder la fin du projet~; le chemin critique est la chaîne des tâches dont la marge est nulle.
\end{definition}

Kelley et Walker ont décrit la méthode en 1959 pour des projets de construction et de maintenance~; elle s'applique sans changement à un \omterm{def:fm:entering-a-new-market:plan}{plan d'entrée}
(le \cref{lst:fm:entering-a-new-market:cpm}).

\omcode{code/firm/entryplan/firm_entryplan.py}{30}{64}{La méthode du chemin critique~: débuts et fins au plus tôt par récurrence, fins au plus tard à rebours,
marges et chemin critique~; et l'échéancier mensuel des coûts.}\label{lst:fm:entering-a-new-market:cpm}

Le chemin critique traverse l'entité juridique, l'enregistrement réglementaire, l'adhésion à la bourse et la certification~: dix mois avant la première
transaction (le \cref{fig:fm:entering-a-new-market:gantt}). La chaîne de compensation, autour de laquelle le plan de huit mois avait été construit, se termine au bout de huit
mois et dispose d'un mois de marge. La connectivité dispose de quatre mois de marge, et les licences de données de six mois~: elles peuvent démarrer tard sans coût. Avant la première
transaction, l'entrée dépense 949\,000~\$, au plus 170\,000~\$ dans un mois, lorsque le logiciel, les embauches, l'enregistrement et la connectivité se chevauchent.

\begin{omfigure}
\begin{tikzpicture}
\begin{axis}[width=10.5cm,height=6.2cm,xbar stacked,bar width=7pt,xmin=0,xmax=11,xlabel={\small mois},
  ytick={0,1,2,3,4,5,6,7,8,9,10},yticklabels={entité juridique,identifiant,enregistrement,convention de compensation,limites de compensation,adhésion,connectivité,
  licences de données,embauches,logiciel,certification},y dir=reverse,ymin=-0.6,ymax=10.6,tick label style={font=\footnotesize},grid=major,
  grid style={gray!20},table/col sep=comma]
\addplot[draw=none,fill=none] table[y=pos,x=start,restrict expr to domain={\coordindex}{0:10}]{figdata/desk/24-entering-a-new-market/gantt.csv};
\addplot[fill=omThm!45,draw=omThm] table[y=pos,x expr=\thisrow{months}*(1-\thisrow{critical}),restrict expr to domain={\coordindex}{0:10}]{figdata/desk/24-entering-a-new-market/gantt.csv};
\addplot[fill=omDef!70,draw=omDef] table[y=pos,x expr=\thisrow{months}*\thisrow{critical},restrict expr to domain={\coordindex}{0:10}]{figdata/desk/24-entering-a-new-market/gantt.csv};
\draw[black,thick] (axis cs:10,-0.6) -- (axis cs:10,10.6) node[pos=0.03,left,font=\scriptsize] {première transaction};
\draw[gray,dashed] (axis cs:8,-0.6) -- (axis cs:8,10.6) node[pos=0.97,left,font=\scriptsize] {plan};
\end{axis}
\end{tikzpicture}

\omcaption{Le calendrier de l'entrée~: chaque chantier à partir de son début au plus tôt, le chemin critique dans la couleur la plus foncée. La première transaction intervient à dix
mois, et non à huit mois comme prévu~: la chaîne de l'enregistrement et de l'adhésion dure deux mois de plus que la chaîne de compensation. Données~:
\texttt{firm.entryplan.schedule}.}\label{fig:fm:entering-a-new-market:gantt}
\end{omfigure}

\section{Échelonner l'investissement}

\begin{definition}[Investissement par étapes]\label{def:fm:entering-a-new-market:staged}
Un \emph{investissement par étapes}\index{investissement par étapes} engage des fonds en tranches séparées par des points de décision, chacune assortie d'un critère d'abandon
fixé à l'avance selon les jalons de validation (stage gates) du livre~7, de sorte que la société puisse s'arrêter lorsque ce qu'elle a appris indique que la suite ne vaut pas la dépense.
\end{definition}

Après la première transaction, le chiffre d'affaires monte vers un régime permanent $R$ par mois, selon $R(1-e^{-t/\tau})$ avec $\tau=4$ mois, face à des coûts d'exploitation de
250\,000~\$ par mois~; l'horizon est de trois ans et le taux d'actualisation de 10~\%. Le régime permanent est la partie incertaine~: dans le modèle du chapitre, il suit
une loi lognormale de médiane 400\,000~\$ par mois. Un scénario faible (150\,000~\$) fait perdre 3.55~millions de dollars de valeur actuelle~; le cas de base (450\,000~\$) rapporte
2.00~millions de dollars, soit un taux de rendement interne de 104~\% par an~; un scénario fort (800\,000~\$) rapporte 8.46~millions de dollars.

\begin{proposition}[Un point d'arrêt ne nuit jamais avec une information parfaite]\label{prop:fm:entering-a-new-market:kill}
Si au point d'arrêt la société apprend exactement $R$ et s'arrête lorsque $R$ est inférieur au coût d'exploitation $c$, la valeur actuelle nette espérée de l'entrée avec le
point d'arrêt est au moins égale à sa valeur sans celui-ci.
\end{proposition}

\begin{proof}
Après le point d'arrêt, poursuivre ajoute $R(1-e^{-t/\tau})-c<R-c<0$ à chaque mois lorsque $R<c$, tandis que s'arrêter n'ajoute rien~; lorsque $R\ge c$, la société poursuit et
rien ne change. Le point d'arrêt augmente donc la valeur dans chaque état où $R<c$ et laisse les autres inchangés~; il en va de même pour son espérance.
\end{proof}

En pratique, la société ne voit qu'une estimation bruitée. La société du chapitre observe $R$ avec une erreur de 25~\% après six mois d'exploitation effective et s'arrête si
l'estimation est inférieure au coût d'exploitation de 250\,000~\$ (\cref{lst:fm:entering-a-new-market:stage}).

\omcode{code/firm/entryplan/firm_entryplan.py}{94}{107}{La valeur de la séquentialité~: pour chaque revenu réel, la VAN sans point d'arrêt et avec un point d'arrêt fondé sur
une estimation bruitée.}\label{lst:fm:entering-a-new-market:stage}

\section{Tutoriel~: huit mois jusqu'à la première transaction}\label{tut:fm:entering-a-new-market:tut}

\begin{tutorial}
\textbf{Objectif.} Tracer le chemin critique de l'entrée, la chiffrer et valoriser un point d'arrêt après six mois d'exploitation. \textbf{État final~:} le diagramme de Gantt
(\cref{fig:fm:entering-a-new-market:gantt}) et les distributions de VAN avec et sans séquentialité (\cref{fig:fm:entering-a-new-market:npv}).
\begin{tutsteps}
\item \textbf{Le réseau.} \texttt{fm\_entry.TASKS}~; \texttt{firm.entryplan.schedule} et \texttt{critical\_path}.
\item \textbf{Les coûts.} \texttt{cost\_by\_month} répartit le coût de chaque chantier sur ses mois.
\item \textbf{Les issues.} \texttt{fm\_entry.scenario\_npvs()} pour trois niveaux de revenu~; \texttt{revenue\_draws} pour 4\,000 tirages lognormaux.
\item \textbf{Le point d'arrêt.} \texttt{fm\_entry.staging()} compare les VAN~; \texttt{threshold\_curve()} fait varier le seuil d'arrêt.
\end{tutsteps}
\end{tutorial}

\begin{omfigure}
\begin{tikzpicture}
\begin{axis}[width=10.5cm,height=5.2cm,ybar,bar width=4pt,xlabel={\small valeur actuelle nette (millions de \$)},ylabel={\small part des issues (\%)},
  xmin=-6.5,xmax=15,ymin=0,ymax=30,tick label style={font=\footnotesize},grid=major,grid style={gray!20},table/col sep=comma,
  legend cell align=left,legend style={font=\scriptsize,at={(0.5,-0.3)},anchor=north,/tikz/every even column/.append style={column sep=8pt}},legend columns=2]
\addplot[fill=gray!45,draw=gray] table[x=mid,y=plain]{figdata/desk/24-entering-a-new-market/npv.csv};
\addplot[fill=omDef!55,draw=omDef] table[x=mid,y=staged]{figdata/desk/24-entering-a-new-market/npv.csv};
\legend{sans point d'arrêt,point d'arrêt après six mois}
\end{axis}
\end{tikzpicture}

\omcaption{La valeur actuelle nette de l'entrée sur 4\,000 tirages de son revenu en régime établi, avec et sans point d'arrêt après six mois d'exploitation.
La séquentialité supprime les pires pertes~: le 5\up{e}~centile passe de $-\$3.57$ millions à $-\$1.92$ millions. Données~: \texttt{fm\_entry.staging}.}\label{fig:fm:entering-a-new-market:npv}
\end{omfigure}

Sans point d'arrêt, la VAN espérée de l'entrée vaut 2.52~millions de dollars, avec une probabilité de 38.6~\% de perdre de l'argent et un 5\up{e}~centile de $-\$3.57$ millions. Avec
le point d'arrêt, la société renonce à 22.7~\% des entrées~; la VAN espérée augmente de 200\,000~\$ pour atteindre 2.72~millions de dollars, et le 5\up{e}~centile remonte à $-\$1.92$
millions. La probabilité de perdre de l'argent bouge à peine (39.6~\%)~: l'arrêt transforme les grosses pertes en pertes plus petites, et non en gains, et les
coûts de mise en place sont irrécupérables. Le bruit a un coût~: 3.6~\% de toutes les entrées sont faibles et se poursuivent, et 5.5~\% sont saines et arrêtées.

\begin{omfigure}
\begin{tikzpicture}
\begin{axis}[width=10.5cm,height=5cm,xlabel={\small seuil d'arrêt (milliers de \$ de revenu mensuel estimé)},ylabel={\small valeur du point d'arrêt (milliers de \$)},
  xmin=0,xmax=700,ymin=-1400,ymax=300,tick label style={font=\footnotesize},grid=major,grid style={gray!20},table/col sep=comma]
\addplot[omDef,thick,mark=*,mark size=1.2pt] table[x=threshold,y=value]{figdata/desk/24-entering-a-new-market/threshold.csv};
\draw[black,thin] (axis cs:0,0) -- (axis cs:700,0);
\draw[gray,dotted] (axis cs:250,-1400) -- (axis cs:250,300) node[pos=0.1,right,font=\scriptsize] {coût d'exploitation};
\end{axis}
\end{tikzpicture}

\omcaption{La valeur que le point d'arrêt ajoute à la VAN espérée de l'entrée, selon le seuil de revenu en dessous duquel la société s'arrête. Elle culmine près du
coût d'exploitation, 250\,000~\$ par mois, et devient négative au-delà de 350\,000~\$, où la règle arrête trop d'entrées saines. Données~:
\texttt{fm\_entry.threshold\_curve}.}\label{fig:fm:entering-a-new-market:threshold}
\end{omfigure}

Le seuil compte autant que le point d'arrêt (la \cref{fig:fm:entering-a-new-market:threshold}). Fixé au niveau du coût d'exploitation, comme
le suggère la \cref{prop:fm:entering-a-new-market:kill}, il ajoute 200\,000~\$~; à 300\,000~\$, il ajoute 150\,000~\$~; à 400\,000~\$, il détruit 138\,000~\$,
car une règle qui exige le scénario de base élimine des entrées qui auraient été rentables. Un critère d'abandon se rédige avant l'entrée, à un niveau que la société pourrait
défendre si le marché se révélait ensuite favorable.

\section{Une entrée pas à pas}

\begin{method}[Planifier une entrée]\label{met:fm:entering-a-new-market:plan}
\begin{enumerate}
\item Énoncer pourquoi l'avantage de la société doit fonctionner sur ce marché, et la taille du gain, avec des éléments probants (normes de recherche du livre~7).
\item Lister chaque chantier avec son responsable, sa durée, son coût et ses prédécesseurs~; tracer le réseau et son chemin critique~; inscrire les responsables des
chantiers critiques à la réunion hebdomadaire.
\item Lancer d'abord les chantiers longs et peu coûteux (enregistrement, compensation)~; lancer les chantiers courts et coûteux (recrutements, connectivité) aussi tard que
leur marge le permet.
\item Rédiger les critères d'abandon avant l'entrée~: quand ils seront mesurés, sur quelles données et à quel seuil.
\item Revoir le plan à chaque jalon de validation (stage gate) au regard du réseau et des critères, et non des dates initiales.
\end{enumerate}
\end{method}

\begin{remark}[Quelle marge utiliser]\label{rem:fm:entering-a-new-market:slack}
La marge est une option sur le retard. Commencer les recrutements deux mois plus tard, dans leur marge, reporte 240\,000~\$ de salaires sans retarder la première
transaction, et conserve la possibilité de ne pas les dépenser si l'enregistrement ou l'accord de compensation échoue d'abord. La même logique vaut pour la connectivité,
avec quatre mois de marge et des frais qui courent à partir de la livraison~: le chantier le moins coûteux à lancer est celui qui fixe la date.
\end{remark}

\section{Construction~: le plan d'entrée}\label{bld:fm:entering-a-new-market:entryplan}

\begin{build}
\bfield{Rôle} Une entrée comme réseau de chantiers~: le chemin critique et les marges, le calendrier des coûts, les rampes de revenu, la VAN et le TRI, et la valeur de la
séquentialité.
\bfield{Interface} \texttt{firm.entryplan}~: \texttt{Task}, \texttt{schedule}, \texttt{critical\_path}, \texttt{cost\_by\_month}, \texttt{npv},
\texttt{irr}, \texttt{entry\_cash}, \texttt{staged\_npv}.
\bfield{Règles} Une tâche commence lorsque tous ses prédécesseurs sont terminés~; le chemin critique a une marge nulle~; une entrée arrêtée interrompt ses coûts d'exploitation et son revenu
au mois de l'arrêt.
\bfield{Tests d'acceptation} \texttt{code/firm/entryplan/tests/}~: le calendrier, la marge et le chemin d'un réseau de quatre tâches~; les coûts, la VAN et le TRI sur des chiffres faits à la main~;
la règle d'arrêt avec une information exacte.
\bfield{Pour aller plus loin} Durées incertaines et probabilité de tenir une date~; coûts issus des modèles de compensation et de frais du livre~1~; plusieurs points d'arrêt.
\end{build}

\begin{omsources}
\item Flow Traders, Annual Report 2025.
\item J.~E.~Kelley et M.~R.~Walker, ``Critical-path planning and scheduling'', Eastern Joint Computer Conference, 1959.
\item A.~K.~Dixit et R.~S.~Pindyck, \emph{Investment under Uncertainty}, 1994.
\end{omsources}

\section{Exercices}

\begin{exercise}[$\star$]\label{exo:fm:entering-a-new-market:1}
Calculer la longueur de la chaîne de compensation et de la chaîne d'enregistrement. Laquelle est critique~?
\end{exercise}

\begin{exercise}[$\star$]\label{exo:fm:entering-a-new-market:2}
Indiquer quels chantiers ont la marge la plus importante et ce qu'elle permet.
\end{exercise}

\begin{exercise}[$\star$]\label{exo:fm:entering-a-new-market:3}
Déterminer le coût de l'entrée avant la première transaction et le mois où elle dépense le plus.
\end{exercise}

\begin{exercise}[$\star\star$]\label{exo:fm:entering-a-new-market:4}
Si l'enregistrement durait deux mois au lieu de quatre, quel serait le \omterm{def:fm:entering-a-new-market:plan}{délai jusqu'à la première transaction}, et quelle chaîne deviendrait critique~?
\end{exercise}

\begin{exercise}[$\star\star$]\label{exo:fm:entering-a-new-market:5}
Démontrer la \cref{prop:fm:entering-a-new-market:kill}, et expliquer pourquoi elle peut échouer avec une estimation bruitée.
\end{exercise}

\begin{exercise}[$\star\star$]\label{exo:fm:entering-a-new-market:6}
Pourquoi le point d'arrêt ne modifie-t-il presque pas la probabilité de perdre de l'argent alors qu'il fait monter le 5\up{e}~centile de 1.6~million de dollars~?
\end{exercise}

\begin{exercise}[$\star\star\star$]\label{exo:fm:entering-a-new-market:7}
\emph{Programmation.} Relancer \texttt{fm\_entry.threshold\_curve()} et trouver le meilleur seuil. Que fait un seuil de 400\,000~dollars~?
\end{exercise}

\begin{exercise}[$\star\star\star$]\label{exo:fm:entering-a-new-market:8}
\emph{Trouver la faille.} «~L'accord de compensation est le goulot d'étranglement~: on le lancera en premier et le reste peut attendre.~»
\end{exercise}

\section{Problème~: huit mois jusqu'à la première transaction}

\begin{problem}[{Problème du week-end --- huit mois jusqu'à la première transaction}]\label{pb:fm:entering-a-new-market:1}
Un responsable du développement commercial doit présenter un \omterm{def:fm:entering-a-new-market:plan}{plan d'entrée} sur un marché de contrats à terme dans une nouvelle région, avec une date et une règle de décision.

\medskip\noindent\textbf{Partie~I --- Pourquoi et quoi.}
\begin{enumerate}
\item Pourquoi une société de trading entrerait-elle sur un nouveau marché~? Donner l'exemple public du chapitre.
\item Définir un \omterm{def:fm:entering-a-new-market:plan}{plan d'entrée} et le \omterm{def:fm:entering-a-new-market:plan}{délai jusqu'à la première transaction}.
\item Lister les chantiers et le chapitre sur lequel chacun s'appuie.
\item Quelles dépendances l'accroche décrit-elle~?
\end{enumerate}

\medskip\noindent\textbf{Partie~II --- Le réseau.}
\begin{enumerate}[resume]
\item Définir la \omterm{def:fm:entering-a-new-market:cpm}{méthode du chemin critique}.
\item Donner le calendrier et le chemin critique.
\item Donner la marge de chaque chantier et ce qu'il faut en faire.
\item Donner le coût avant la première transaction et son mois de pointe.
\end{enumerate}

\medskip\noindent\textbf{Partie~III --- L'argent.}
\begin{enumerate}[resume]
\item Décrire la montée en charge des revenus et les VAN des trois scénarios.
\item Définir un \omterm{def:fm:entering-a-new-market:staged}{investissement par étapes}.
\item Énoncer et démontrer la \cref{prop:fm:entering-a-new-market:kill}.
\item Donner la distribution des VAN avec et sans point d'arrêt.
\item Comment faut-il fixer le seuil d'arrêt~?
\end{enumerate}

\medskip\noindent\textbf{Partie~IV --- Le plan.}
\begin{enumerate}[resume]
\item Quels chantiers faudrait-il lancer en premier, et lesquels en dernier~?
\item Qui doit assister à la réunion hebdomadaire~?
\item Que faudrait-il rapporter à chaque jalon de validation (stage gate)~?
\item Qu'est-ce qui conduirait à raccourcir le point d'arrêt à trois mois~?
\item Que laisse de côté le modèle~?
\item Énoncer le \emph{résultat nommé}~: le \omterm{def:fm:entering-a-new-market:plan}{délai jusqu'à la première transaction} sur le chemin critique, et la valeur que le point d'arrêt ajoute à la VAN de l'entrée.
\item En deux phrases, rédiger le plan.
\end{enumerate}
\end{problem}

\section{Questions d'entretien}

\begin{interviewq}[$\star$ \iqroles{développeur}]\label{iq:fm:entering-a-new-market:1}
Calculez le chemin critique d'un petit projet à partir de ses tâches et de ses dépendances.
\end{interviewq}

\begin{interviewq}[$\star$ \iqroles{trader}]\label{iq:fm:entering-a-new-market:2}
Qu'est-ce qui doit être en place avant qu'une société puisse négocier des contrats à terme sur une nouvelle bourse~?
\end{interviewq}

\begin{interviewq}[$\star\star$ \iqroles{chercheur}]\label{iq:fm:entering-a-new-market:3}
Comment évalueriez-vous l'option d'abandonner une nouvelle activité après six mois~?
\end{interviewq}

\begin{interviewq}[$\star\star$ \iqroles{risque}]\label{iq:fm:entering-a-new-market:4}
Quels critères d'abandon fixeriez-vous pour l'entrée sur un nouveau marché, et pourquoi les fixer avant d'entrer~?
\end{interviewq}

\begin{interviewq}[$\star\star$ \iqroles{trader, chercheur}]\label{iq:fm:entering-a-new-market:5}
Pourquoi une société rentable ailleurs pourrait-elle perdre de l'argent en entrant dans une nouvelle région~?
\end{interviewq}

\begin{interviewq}[$\star\star\star$ \iqroles{développeur, chercheur}]\label{iq:fm:entering-a-new-market:6}
Les durées sont incertaines. Comment estimeriez-vous la probabilité de réaliser sa première transaction avant une date donnée~?
\end{interviewq}
